\begin{Exercise}[label={ex:somma_conv}]
    Convertire in base 2 i numeri 77 e 156. Sommare in base 2 i risultati ottenuti.
\end{Exercise}

\begin{Answer}[ref={ex:somma_conv}]
    Per convertire il valore $77_{10}$ saranno necessari $n \geq \lceil \log_2(77) \rceil = \lceil 6.26 \rceil = 7$ bit. 

    \begin{center}
    \begin{tabular}{r|r}
        \multicolumn{2}{c}{2} \\
        77 & 1 \\
        38 & 0 \\
        19 & 1 \\
        9 & 1 \\
        4 & 0 \\
        2 & 0 \\
        1 & 1 \\
        0 &
    \end{tabular}
    \end{center}

    $77_{10} = 1001101_2$
    
    Per convertire il valore $156_{10}$ saranno necessari $n \geq \lceil \log_2(156) \rceil = \lceil 7.28 \rceil = 8$ bit. 

    \begin{center}
    \begin{tabular}{r|r}
        \multicolumn{2}{c}{2} \\
        156 & 0 \\
        78 & 0 \\
        39 & 1 \\
        19 & 1 \\
        9 & 1 \\
        4 & 0 \\
        2 & 0 \\
        1 & 1 \\
        0 &
    \end{tabular}
    \end{center}

    $156_{10} = 10011100_2$
    
    Eseguiamo ora la somma binaria:
    \begin{table}[H]
    \centering
    \begin{tabular}{ccccccccccc}
            &      & 0 & 0 & 1 & 1 & 1 & 1 & 0 & 0 & \\
            &      & 0 & 1 & 0 & 0 & 1 & 1 & 0 & 1 &+\\
            &      & 1 & 0 & 0 & 1 & 1 & 1 & 0 & 0 &=\\
            \hline
            & & 1 & 1 & 1 & 0 & 1 & 0 & 0 & 1 &\\
    \end{tabular}
    \end{table}
    Il risultato è $S = (11101001)_{2}$.
\end{Answer}
